\chapter{Engenharia de Software}
\label{cap:eng-software-teoria}

\begin{edital}
Engenharia de requisitos (classificação, processo, técnicas de elicitação);
testes (unitários, integração, ágeis, usabilidade, automatizados, TDD, ciclo
de vida, RPA); metodologias ágeis (Scrum, Kanban, XP); padrões de
desenvolvimento e reuso; codificação; Ponto de Função e Story Points; DevOps;
design de software.
\end{edital}

Engenharia de software é a disciplina de organizar \textbf{como} um sistema é
construído, não o que ele faz por dentro (isso é o Capítulo
\ref{cap:programacao-teoria}) nem como seus componentes se distribuem (isso é
o Capítulo \ref{cap:arquitetura-teoria}). Este capítulo segue uma ordem
lógica: primeiro os \textbf{modelos de processo} (o “formato” do projeto do
início ao fim), depois \textbf{maturidade organizacional}, depois os
\textbf{blocos que preenchem qualquer processo} --- requisitos, testes,
metodologias ágeis, medição, entrega e design.

\section{Modelos de ciclo de vida}

\begin{conceito}[O que um modelo de ciclo de vida decide]
Todo projeto de software passa pelas mesmas atividades essenciais:
levantar requisitos, projetar, construir, testar, implantar. O que muda de
um \textbf{modelo de ciclo de vida} para outro não é \textbf{o que} se faz,
é \textbf{a ordem e a repetição}: as atividades acontecem uma vez, em
sequência rígida? Repetem em ciclos? Entregam aos pedaços?

\begin{itemize}
\item \textbf{Cascata (\textit{waterfall})}: fases \textbf{sequenciais e
estanques} --- uma termina, a próxima começa, sem voltar atrás.
Requisitos \textbf{congelados} logo no início. Funciona quando o problema é
\textbf{bem compreendido e estável} (ex.: sistema regulatório com
especificação fechada por lei). Falha quando o requisito muda no meio do
caminho, porque o modelo não tem mecanismo para incorporar mudança sem
recomeçar.
\item \textbf{Incremental}: o sistema é entregue em \textbf{pedaços
funcionais que se somam} ao longo do tempo --- cada incremento já é
utilizável. Reduz o risco de “só ver o produto pronto no fim”.
\item \textbf{Iterativo}: o mesmo produto é \textbf{refinado} em ciclos
sucessivos (a primeira versão é rascunho grosseiro, a segunda melhora, e
assim por diante) --- o foco é refinar, não necessariamente entregar pedaço
novo a cada ciclo (a diferença com incremental é sutil e por isso a maioria
dos processos modernos combina os dois: “iterativo-incremental”).
\item \textbf{Espiral (Boehm)}: iterativo, mas com um passo extra em cada
volta da espiral --- \textbf{análise de risco explícita} antes de seguir
adiante. Nasceu para projetos grandes e arriscados, onde o custo de errar é
alto o bastante para justificar parar e avaliar risco a cada ciclo.
\item \textbf{Prototipação}: constrói uma versão \textbf{parcial e rápida}
(às vezes descartável) só para \textbf{validar requisito com o cliente}
antes de investir no sistema real --- serve quando o requisito é
\textbf{incerto ou volátil}, o oposto do cenário ideal para cascata.
\item \textbf{Ágil}: iterativo-incremental com ciclos \textbf{curtos} e
\textbf{adaptação contínua} --- assume que requisito muda e projeta o
processo em torno disso, em vez de tentar blindar o processo contra
mudança.
\end{itemize}
\end{conceito}

\begin{exemplo}[Reconhecendo o modelo pelo enunciado]
“O projeto define todo o escopo antes de iniciar a codificação; cada fase
só começa quando a anterior está formalmente aprovada” $\to$ \textbf{cascata}
(sequencial, sem sobreposição, aprovação formal entre fases). “A equipe
constrói uma versão simplificada da tela principal só para o cliente
confirmar o fluxo antes de programar o sistema de verdade” $\to$
\textbf{prototipação} (o objetivo é validar requisito, não entregar
produto).
\end{exemplo}

\begin{cuidado}[Incremental e iterativo não são sinônimos, embora andem juntos]
\textbf{Incremental} responde “o que é entregue” (pedaços que se somam,
cada um com valor de uso). \textbf{Iterativo} responde “como se chega lá”
(ciclos que refinam o que já existe). É possível ser incremental sem ser
muito iterativo (entregar módulo A, depois B, depois C, sem voltar a
refinar A) e vice-versa. Scrum, na prática, é os dois ao mesmo tempo: cada
Sprint refina o produto (iterativo) e normalmente adiciona um incremento
utilizável (incremental) --- por isso os processos ágeis modernos costumam
ser descritos com o termo composto.
\end{cuidado}

\subsection{Modelo V: cada fase de construção tem seu par de teste}

\begin{conceito}[A lógica “quem definiu, confere”]
O Modelo V é uma variação do cascata desenhada em V: o \textbf{ramo
descendente} (esquerda) vai do requisito ao código, detalhando cada vez
mais; o \textbf{ramo ascendente} (direita) sobe testando, e cada nível de
teste \textbf{verifica exatamente o que foi decidido no nível espelhado do
outro lado}.

\begin{center}
\begin{tabular}{@{}p{6cm}p{5.6cm}@{}}
\toprule
\textbf{Fase (ramo descendente)} & \textbf{Nível de teste que a verifica} \\
\midrule
Requisitos do usuário & \textbf{Teste de aceitação} \\
Especificação do sistema & Teste de sistema \\
Projeto arquitetural (módulos e interfaces) & \textbf{Teste de integração} \\
Projeto detalhado (algoritmos) & Teste de unidade \\
Codificação & (vértice do V) \\
\bottomrule
\end{tabular}
\end{center}

Por que essa lógica faz sentido: quem \textbf{acordou} o requisito com o
usuário é quem consegue dizer, no fim, se ele foi atendido --- daí requisito
$\leftrightarrow$ aceitação. Quem \textbf{decidiu} a divisão em módulos e
suas interfaces (projeto arquitetural) é quem consegue verificar se os
módulos, integrados, funcionam como planejado --- daí arquitetura
$\leftrightarrow$ integração. O par nunca cruza: requisito não casa com
teste de unidade (unidade verifica o \textbf{projeto detalhado}, não a
visão do usuário), e código não casa com aceitação.
\end{conceito}

\subsection{Verificação x validação}

\begin{conceito}[Duas perguntas diferentes sobre o mesmo software]
O Modelo V chama de “verificação” todo o ramo ascendente, mas os dois
termos não são sinônimos --- e a confusão entre eles é o distrator mais
recorrente de todo o bloco, porque as duas palavras soam parecidas em
português e a resposta muda inteira conforme qual delas o enunciado
descreve.

\begin{center}
\begin{tabular}{@{}p{2.6cm}p{4.4cm}p{4.6cm}@{}}
\toprule
& \textbf{Verificação} & \textbf{Validação} \\
\midrule
A pergunta & “estamos construindo o produto \textbf{corretamente}?” &
“estamos construindo o \textbf{produto certo}?” \\
Compara com & a \textbf{especificação} (documento da fase anterior) &
a \textbf{necessidade real} do usuário \\
Quando & ao longo de todo o desenvolvimento, fase a fase & principalmente
perto da entrega, com o usuário \\
Como se faz & revisão, inspeção, \textit{walkthrough}, análise estática,
teste de unidade e integração & teste de aceitação, teste alfa/beta,
homologação, protótipo avaliado pelo cliente \\
Precisa executar o software? & \textbf{não necessariamente} --- revisar
documento já é verificar & na prática, \textbf{sim} \\
\bottomrule
\end{tabular}
\end{center}

\textbf{A âncora de uma palavra:} verificação olha para o \textit{papel} (a
especificação); validação olha para a \textit{pessoa} (o usuário).

A consequência prática que resolve qualquer cenário: \textbf{um sistema pode
passar na verificação e falhar na validação}. Ele implementa exatamente o
que o documento pedia --- só que o documento pedia a coisa errada. É o
sistema que atende a cada linha do requisito e não serve para o trabalho de
ninguém: verificação impecável, validação reprovada. O caminho inverso
também existe, mais raro: o produto resolve a necessidade real apesar de
divergir do que a especificação escrita dizia.
\end{conceito}

\begin{cuidado}[Os dois erros de troca mais comuns]
Definir \textbf{validação} como “conformidade com a especificação” (isso é
verificação); ou definir \textbf{verificação} como “atendimento às
necessidades do usuário” (isso é validação). Um terceiro erro: tratar as
duas como sinônimos de “teste” --- teste é uma \textbf{técnica}, usada
pelas duas (existe teste de unidade, que verifica, e teste de aceitação,
que valida). No Modelo V, a ponta de baixo (unidade/integração/sistema) é
verificação; a ponta de cima (aceitação) é validação.
\end{cuidado}

\subsection{RUP: fase x disciplina}

\begin{conceito}[Os dois eixos independentes do RUP]
O RUP (\textit{Rational Unified Process}) é iterativo-incremental e dirigido
por casos de uso. A estrutura que sustenta qualquer questão sobre ele tem
dois eixos que \textbf{não se misturam}:

\textbf{Fase} é o eixo do \textbf{tempo} --- o projeto passa por Concepção,
Elaboração, Construção e Transição, \textbf{nessa ordem, uma única vez},
cada fase terminando num marco de decisão. \textbf{Disciplina}
(\textit{workflow}) é o eixo do \textbf{conteúdo} --- Requisitos, Análise e
Projeto, Implementação, Teste, Implantação, mais as de apoio (Gerência de
Projeto, Configuração e Mudança, Ambiente).

O ponto central: \textbf{as disciplinas atravessam todas as fases}. Não
existe “a fase de testes” --- testa-se em todas as quatro fases, só que em
intensidade diferente a cada uma (o gráfico clássico do RUP desenha as
fases como colunas e as disciplinas como linhas horizontais; cada linha tem
uma “corcova” de maior intensidade em certa coluna, mas nenhuma linha
começa e termina dentro de uma coluna só).

\begin{center}
\begin{tabular}{@{}p{2.4cm}p{5.2cm}p{4.4cm}@{}}
\toprule
\textbf{Fase} & \textbf{Pergunta que ela responde} & \textbf{Marco} \\
\midrule
Concepção & vale a pena fazer? escopo, viabilidade, caso de negócio & Objetivos do ciclo de vida \\
\textbf{Elaboração} & \textbf{como fazer sem que dê errado?} --- ataca os \textbf{riscos} e firma a \textbf{arquitetura executável} & Arquitetura do ciclo de vida \\
Construção & constrói o grosso do produto, iteração a iteração & Capacidade operacional inicial \\
Transição & entrega ao usuário: implantação, treinamento, ajuste & Liberação do produto \\
\bottomrule
\end{tabular}
\end{center}

\textbf{Por que risco e arquitetura vivem na Elaboração:} a lógica é
econômica --- resolver o risco arquitetural \textbf{antes} de gastar o
orçamento da Construção. Descobrir um problema de arquitetura já na
Construção é caro demais para corrigir.
\end{conceito}

\begin{cuidado}[Erros de leitura comuns no RUP]
Tratar disciplina como fase (“a fase de testes do RUP” não existe: teste é
disciplina, presente nas quatro fases). Deslocar o risco para a Concepção
(lá se avalia viabilidade, não se ataca risco técnico) ou para a Construção
(lá já é tarde e caro). Chamar o RUP de “cascata”: ele é iterativo
\textbf{dentro} de cada fase; o que é sequencial é só a \textbf{ordem} das
quatro fases, não o trabalho interno de cada uma.
\end{cuidado}

\section{Maturidade de processo: CMMI e MPS.BR}

\begin{conceito}[Que problema a maturidade de processo resolve]
Modelos de maturidade avaliam \textbf{quão maduro é o processo} de uma
organização --- não a qualidade de um produto específico. O CMMI nasceu de
uma pergunta prática do Departamento de Defesa americano nos anos 1980:
\textit{como saber, antes de contratar, se um fornecedor vai entregar?}
Olhar o produto anterior não bastava --- um projeto que deu certo pode ter
dado certo por acaso, ou porque alguém virou noites. O que se queria medir
era a \textbf{capacidade de repetir} o resultado.

Daí a definição central: \textbf{maturidade é o quanto o resultado depende
do processo, e não das pessoas}. Uma organização madura entrega parecido
com equipes diferentes, porque o jeito de trabalhar está descrito, é
seguido e é medido. Uma organização imatura entrega bem quando calha de ter
a pessoa certa no projeto certo.
\end{conceito}

\begin{conceito}[CMMI --- os cinco níveis, por estágios]
\begin{center}
\begin{tabular}{@{}p{1cm}p{3.6cm}p{7.2cm}@{}}
\toprule
\textbf{Nº} & \textbf{Nível} & \textbf{O que caracteriza} \\
\midrule
1 & Inicial & processo imprevisível, reativo, dependente de heróis \\
2 & Gerenciado & gerenciado \textbf{por projeto} (cada projeto planeja e
monitora do seu jeito) \\
3 & \textbf{Definido} & processo \textbf{padronizado no âmbito da
organização}, não projeto a projeto \\
4 & Gerenciado Quantitativamente & processo \textbf{medido e controlado por
estatística} \\
5 & Em Otimização & melhoria contínua a partir da medição \\
\bottomrule
\end{tabular}
\end{center}

O corte mais importante para entender (e o mais cobrado) é do nível 2 para
o 3: no \textbf{nível 2}, o projeto A tem cronograma e controle de mudanças,
o projeto B ao lado faz diferente --- e os dois podem estar bem gerenciados,
porque a disciplina existe \textbf{por projeto}. No \textbf{nível 3}, existe
um \textbf{processo-padrão da organização}, e cada projeto o adapta
(\textit{tailoring}) dentro de regras --- um desenvolvedor que muda de
projeto reconhece o jeito de trabalhar, e o aprendizado de um projeto vira
ativo de todos.

Do 3 para cima, a evolução é sobre \textbf{medir}: no \textbf{nível 4}, a
organização não diz mais “costuma dar certo”, diz “a densidade de defeito
fica entre tais limites” --- previsibilidade quantitativa. No \textbf{nível
5}, essa medição é usada para \textbf{mudar o processo de propósito} e
verificar se a mudança realmente melhorou o resultado --- melhoria contínua
baseada em dado, não em opinião.

\textbf{Duas representações, duas perguntas diferentes:} a representação
\textbf{por estágios} (a tabela acima) responde “que nota tem esta
\textbf{organização}?” --- um número só, comparável entre fornecedores, e
foi para isso que o modelo nasceu. A representação \textbf{contínua}
responde “quão \textbf{capaz} somos \textbf{nesta área específica}
(por exemplo, só Engenharia de Requisitos)?” --- permite investir numa área
sem ter de elevar tudo junto. Por isso o vocabulário muda: estágios dá
\textbf{maturidade} à organização inteira; contínua dá \textbf{capacidade} a
cada área de processo isolada.
\end{conceito}

\begin{conceito}[MPS.BR --- a versão brasileira, de G a A]
Mesma ideia do CMMI, escala própria (MR-MPS-SW) com \textbf{sete níveis
identificados por letras}, evoluindo \textbf{de G para A}: G (Parcialmente
Gerenciado), F (Gerenciado), E (Parcialmente Definido), D (Largamente
Definido), C (Definido), B (Gerenciado Quantitativamente), A (Em
Otimização). O G é o ponto de partida (organização menos madura); o A é o
topo.
\end{conceito}

\begin{cuidado}[Trocas clássicas nos dois modelos]
Chamar o nível 3 do CMMI de “Gerenciado” (ele é o \textbf{Definido} --- o
Gerenciado é o 2). Inverter a ordem do topo, pondo Gerenciado
Quantitativamente depois de Em Otimização (a ordem certa é Definido $\to$
Gerenciado Quantitativamente $\to$ Em Otimização). Achar que o MPS.BR
começa no A --- é o contrário, começa no G. E lembrar que a escala tem
\textbf{cinco} níveis no CMMI e \textbf{sete} no MPS.BR --- não são a mesma
contagem.
\end{cuidado}

\section{Engenharia de requisitos}

\subsection{Requisito funcional x não funcional}

\begin{conceito}[O que o sistema faz x quão bem ele faz]
\begin{center}
\begin{tabular}{@{}p{2.6cm}p{5.4cm}p{5cm}@{}}
\toprule
\textbf{Tipo} & \textbf{O que descreve} & \textbf{Exemplo} \\
\midrule
\textbf{Funcional (RF)} & O QUE o sistema faz --- uma função, um
comportamento observável & “emitir saldo da conta” \\
\textbf{Não funcional (RNF)} & COMO / QUÃO BEM --- uma qualidade ou
restrição sobre como as funções se comportam & “saldo em tempo real”,
desempenho, segurança, usabilidade, disponibilidade \\
\bottomrule
\end{tabular}
\end{center}

O mecanismo por trás da distinção: um RF sobrevive sozinho como frase com
verbo de ação (“o sistema deve emitir”, “o usuário deve poder cadastrar”).
Um RNF \textbf{qualifica} um RF ou o sistema como um todo --- ele não tem
sentido sem alguma função para qualificar (“rápido” --- rápido fazendo o
quê?). Por isso RNF costuma vir como adjetivo/advérbio de qualidade
(“seguro”, “disponível 99,9\% do tempo”, “em tempo real”), enquanto RF vem
como verbo de ação.
\end{conceito}

\begin{exemplo}[Separando RF de RNF num enunciado]
“O sistema deve permitir que o usuário emita boletos (1). O sistema deve
processar até 1000 requisições por segundo (2). O boleto deve ser emitido
em até 2 segundos após a solicitação (3). O sistema deve estar em
conformidade com a LGPD (4).” --- (1) é RF (ação: emitir boletos); (2), (3)
e (4) são RNF (desempenho, desempenho/tempo de resposta, e conformidade
legal, respectivamente --- este último é um RNF do tipo “externo/legal”).
\end{exemplo}

\begin{regrapratica}
Ao ler um enunciado longo, marque cada \textbf{verbo de ação} do sistema
(RF) e cada palavra de \textbf{qualidade/restrição} (“deve ser
rápido/seguro/disponível/acessível/auditável”) (RNF). RNF se subdivide em
tipos (desempenho, usabilidade, segurança, produto, organizacional,
externo/legal) --- quando a questão pede o \textbf{tipo} de RNF, não só
“é RNF”, associe a palavra-chave ao tipo correspondente.
\end{regrapratica}

\subsection{Processo de engenharia de requisitos}

\begin{conceito}[As cinco etapas, em ordem]
\textbf{Elicitação} (levantamento) $\to$ \textbf{análise} (entender,
priorizar, resolver conflitos) $\to$ \textbf{especificação} (documentar
formalmente) $\to$ \textbf{validação} (confirmar com o interessado que é
isso mesmo) $\to$ \textbf{gerenciamento} (rastrear mudanças ao longo do
projeto --- os requisitos evoluem, e cada mudança precisa de rastreabilidade
até o que ela impacta).

Técnicas de \textbf{elicitação} (todas legítimas, escolhidas conforme o
contexto): entrevista, questionário, \textbf{brainstorming}, workshop,
observação (o analista acompanha o usuário no trabalho real, útil quando o
usuário não sabe verbalizar o que faz), prototipação, análise de documentos
existentes (manuais, sistemas legados).
\end{conceito}

\begin{cuidado}
Brainstorming não é técnica “informal demais para prova de requisito” ---
é técnica de elicitação plenamente válida e cobrada como tal. E a ordem do
processo é sempre elicitar \textbf{antes} de codificar, nunca o inverso ---
requisito que “surge depois da implementação” é sinal de processo mal
seguido, não de uma etapa válida do processo.
\end{cuidado}

\section{Metodologias ágeis}

\subsection{Scrum: framework, não metodologia}

\begin{conceito}[Por que “framework” e não “metodologia”]
Scrum define \textbf{papéis, eventos e artefatos} --- uma estrutura mínima
--- mas não dita \textbf{como} o time deve executar cada atividade dentro
dela (isso é “metodologia”: um passo a passo prescritivo). O Scrum Guide é
propositalmente enxuto: ele diz “existe um Sprint Planning”, não “faça o
Sprint Planning exatamente assim”. Por isso é chamado de framework: a
moldura é fixa, o conteúdo dentro dela é decidido pelo time.

O trabalho é organizado em \textbf{Sprints}: ciclos de tamanho fixo
(\textit{time-boxed}), normalmente de 1 a 4 semanas, sempre com a mesma
duração dentro de um mesmo projeto.

\textbf{Papéis (\textit{accountabilities}, termo do Scrum Guide 2020):}
\begin{itemize}
\item \textbf{Product Owner}: dono do \textbf{valor} do produto --- decide
o que entra no Product Backlog e em que ordem, representando o
interesse do negócio/cliente.
\item \textbf{Scrum Master}: \textbf{remove impedimentos} e serve ao time
--- garante que o Scrum seja entendido e praticado, mas \textbf{não manda}
no time nem distribui tarefas.
\item \textbf{Developers}: quem constrói o incremento a cada Sprint.
\end{itemize}

\textbf{Eventos:} a Sprint em si, mais quatro eventos dentro dela ---
\textbf{Sprint Planning} (o que será feito e como, no início), \textbf{Daily
Scrum} (sincronização diária de 15 minutos), \textbf{Sprint Review}
(inspeciona o \textbf{incremento} com os interessados, no fim), \textbf{Sprint
Retrospective} (inspeciona o \textbf{processo} do time, também no fim, após
a Review).

\textbf{Artefatos e seus compromissos (Scrum Guide 2020):} cada artefato tem
um compromisso associado, que existe para dar foco e permitir medir
progresso --- três pares fixos que não se misturam:

\begin{center}
\begin{tabular}{@{}p{4.2cm}p{4.4cm}p{4cm}@{}}
\toprule
\textbf{Artefato} & \textbf{Compromisso} & \textbf{O que ele fixa} \\
\midrule
Product Backlog & \textbf{Meta do Produto} (Product Goal) & o objetivo de
longo prazo \\
Sprint Backlog & \textbf{Meta da Sprint} (Sprint Goal) & o objetivo único
da Sprint \\
Increment & \textbf{Definition of Done} & o padrão de qualidade do
“pronto” \\
\bottomrule
\end{tabular}
\end{center}
\end{conceito}

\begin{exemplo}[Review x Retrospective num mesmo dia]
No último dia de uma Sprint de duas semanas, o time faz duas reuniões
seguidas. Na primeira, mostra a nova funcionalidade de exportar relatórios
para o cliente e outros interessados, que dão feedback sobre o que
priorizar a seguir --- isso é \textbf{Sprint Review} (foco no
\textbf{produto}). Na segunda, só o time se reúne e discute que o
retrabalho de testes atrasou a Sprint, e decide mudar a forma de escrever
casos de teste na próxima --- isso é \textbf{Sprint Retrospective} (foco no
\textbf{processo}).
\end{exemplo}

\begin{cuidado}[Ligações que a confusão costuma trocar]
\textbf{Sprint Review = produto}; \textbf{Sprint Retrospective = processo}
--- inverter as duas é o erro mais comum. Nos compromissos, a confusão
mais comum é ligar Definition of Done ao Product Backlog ou ao Sprint
Backlog (ela é sempre do \textbf{Incremento}) e a Meta do Produto ao
Incremento (ela é sempre do \textbf{Product Backlog}). Guie-se pelo
horizonte de tempo: produto (longo prazo) $\to$ Product Backlog; Sprint
(curto prazo, uma iteração) $\to$ Sprint Backlog; “está pronto ou não”
$\to$ Incremento. E no Daily, quando surge risco à Sprint, o Scrum Master
\textbf{facilita a solução do time} --- não redistribui tarefas sozinho,
nem assume a tarefa do desenvolvedor: quem decide como resolver é o
próprio time.
\end{cuidado}

\subsection{Kanban, XP, Lean}

\begin{conceito}
\begin{center}
\begin{tabular}{@{}p{2.4cm}p{9.8cm}@{}}
\toprule
\textbf{Framework} & \textbf{Ideia central} \\
\midrule
Kanban & Fluxo \textbf{contínuo}, sem ciclos fechados (sem Sprint); limita o
\textbf{WIP} (\textit{work in progress}, trabalho simultâneo em curso) para
evitar sobrecarga; entrega assim que cada item termina, não em lote. \\
XP (\textit{Extreme Programming}) & Conjunto de \textbf{práticas de
engenharia}: programação em par (\textit{pair programming}), TDD,
integração contínua, refatoração constante, cliente presente no time. \\
Lean & Elimina desperdício e otimiza o fluxo de valor de ponta a ponta ---
originado na manufatura (Toyota), adaptado ao software. \\
Ágil híbrido & Combina práticas ágeis com elementos de métodos
tradicionais (por exemplo, um cronograma macro fixo por contrato, com
execução interna em Sprints). \\
\bottomrule
\end{tabular}
\end{center}
\end{conceito}

\begin{cuidado}
Kanban \textbf{não} tem Sprint --- “Kanban organiza o trabalho em sprints”
é sempre falso. Scrum \textbf{tem} \textit{time-box} fixo; Kanban é fluxo
contínuo, sem prazo fechado por ciclo (o controle é pelo limite de WIP, não
pelo calendário).
\end{cuidado}

\section{Testes de software}

\subsection{Níveis e tipos}

\begin{conceito}
\begin{center}
\begin{tabular}{@{}p{3cm}p{9.2cm}@{}}
\toprule
\textbf{Nível/tipo} & \textbf{O que verifica} \\
\midrule
Unitário & a menor unidade isolada (função/método), sem depender de outras
partes do sistema \\
Integração & a interação entre módulos/componentes já unidos \\
Sistema & o sistema completo, de ponta a ponta \\
Aceitação & se o sistema atende ao que o cliente/usuário precisava (é o
teste que \textbf{valida}, na terminologia da Seção anterior) \\
Usabilidade & se a interface é intuitiva e a experiência de uso é boa \\
Regressão & se uma mudança recente não quebrou algo que já funcionava \\
Fumaça (\textit{smoke}) & verificação rápida e superficial das funções
principais, para decidir se vale a pena testar o resto \\
\bottomrule
\end{tabular}
\end{center}
\end{conceito}

\subsection{Testes de desempenho: carga, estresse, volume}

\begin{conceito}[Os três se confundem porque todos “sobrecarregam” o sistema]
O que muda entre eles é \textbf{até onde} se vai e \textbf{o que} está
sendo sobrecarregado:

\begin{itemize}
\item \textbf{Carga (\textit{load})}: roda o volume de uso \textbf{esperado
em produção} e observa se o sistema aguenta o que foi previsto/contratado.
\item \textbf{Estresse (\textit{stress})}: vai \textbf{além} do previsto,
aumentando progressivamente até achar o \textbf{ponto de ruptura} --- e
observa como o sistema se comporta ao quebrar (degrada com elegância? cai
tudo de uma vez? se recupera sozinho depois?).
\item \textbf{Volume}: sobrecarrega com muita \textbf{massa de dados
armazenada} (um banco com 500 milhões de registros), não necessariamente
muitos usuários simultâneos.
\end{itemize}
\end{conceito}

\begin{cuidado}
“Ultrapassar o previsto até o sistema deixar de responder” é
\textbf{estresse}; “comportar-se bem no volume contratado/esperado” é
\textbf{carga}. O gatilho textual é a palavra \textbf{além}/\textbf{ultrapassa}
(estresse) versus \textbf{esperado}/\textbf{previsto} (carga). Volume é
massa de \textbf{dados}, não de usuários simultâneos --- não confunda com
carga só porque os dois envolvem “muito”.
\end{cuidado}

\subsection{Técnicas de teste caixa-preta: partição de equivalência e valor limite}

\begin{conceito}[Duas técnicas para escolher casos de teste sem testar tudo]
Testar \textbf{toda} entrada possível é inviável (um campo numérico de 0 a
1000 tem 1001 valores possíveis). As duas técnicas clássicas de caixa-preta
resolvem isso escolhendo \textbf{poucos} casos que representam bem o espaço
de entrada:

\begin{itemize}
\item \textbf{Partição de equivalência}: divide as entradas possíveis em
\textbf{classes} que, supõe-se, o sistema trata do mesmo jeito --- testando
\textbf{um} valor de cada classe, testa-se a classe inteira. Para um campo
que aceita idade de 18 a 65, há três classes: “menor que 18” (inválida),
“18 a 65” (válida), “maior que 65” (inválida) --- basta um valor de cada.
\item \textbf{Análise de valor limite}: em vez de um valor qualquer da
classe, testa \textbf{exatamente as fronteiras} entre classes (17, 18, 65,
66) --- porque é ali que a maioria dos defeitos reais de código acontece
(erro de “$>$” no lugar de “$\geq$”, por exemplo).
\end{itemize}
As duas se complementam: partição de equivalência decide \textbf{onde}
procurar; valor limite decide o \textbf{ponto exato} mais eficiente dentro
de cada fronteira.
\end{conceito}

\subsection{Cobertura de caixa-branca}

\begin{conceito}[Critérios que medem quanto do código os testes exercitam]
\begin{itemize}
\item \textbf{Cobertura de comandos} (\textit{statement}): cada linha
executada ao menos uma vez.
\item \textbf{Cobertura de decisões} (\textit{branch}): cada decisão
avaliada \textbf{como verdadeira e como falsa} pelo menos uma vez.
\item \textbf{Cobertura de caminhos}: cada combinação possível de
caminhos pelo código --- o critério mais forte, e em geral inviável de
atingir por completo (o número de caminhos cresce exponencialmente com o
número de decisões).
\end{itemize}
Força crescente: comandos $<$ decisões $<$ caminhos. \textbf{100\% de
cobertura de decisões implica 100\% de cobertura de comandos; a recíproca é
falsa} --- cobrir toda linha não garante ter testado toda combinação de
verdadeiro/falso de cada condição.
\end{conceito}

\begin{exemplo}[Por que 100\% de comandos não é 100\% de decisões]
Código: \texttt{if (a > 10) \{ x = 1; \}} --- um \texttt{if} sem
\texttt{else}. Um único caso de teste com \texttt{a = 20} faz \textbf{todas
as linhas} executarem: cobertura de comandos em 100\%. Mas a condição
\texttt{a > 10} nunca foi avaliada como \textbf{falsa} nesse teste ---
cobertura de decisões fica em 50\%, mesmo sem existir nenhum bloco
\texttt{else} escrito no código. Um segundo caso com \texttt{a = 5}
completaria a cobertura de decisões.
\end{exemplo}

\subsection{TDD, BDD, caixa-preta x caixa-branca, RPA}

\begin{conceito}
\textbf{TDD (\textit{Test-Driven Development})}: escreve o \textbf{teste
antes} do código de produção, no ciclo \textbf{red} (teste falha, porque o
código ainda não existe) $\to$ \textbf{green} (escreve o código mínimo que
faz o teste passar) $\to$ \textbf{refactor} (melhora o código mantendo os
testes passando). O foco é tanto garantir cobertura quanto \textbf{guiar o
design} --- escrever o teste primeiro obriga a pensar na interface antes da
implementação.

\textbf{BDD (\textit{Behavior-Driven Development})}: evolução do TDD que
especifica o comportamento esperado em linguagem \textbf{acessível ao
negócio}, não só a quem programa --- tipicamente em Gherkin
(\textit{Dado}/\textit{Quando}/\textit{Então}), permitindo que
stakeholders não técnicos leiam e validem os cenários de teste.

\textbf{Caixa-preta} (testa entrada/saída, sem olhar o código por dentro) x
\textbf{caixa-branca} (baseado na estrutura interna/no código-fonte) ---
partição de equivalência e valor limite são caixa-preta; cobertura de
comandos/decisões/caminhos é caixa-branca.

\textbf{RPA (\textit{Robotic Process Automation})}: automatiza tarefas
repetitivas de \textbf{interface} (robôs de software que clicam, digitam,
copiam dados entre sistemas) --- \textbf{não é técnica de teste}, é
automação de processo de negócio; aparece no edital de eng.\ de software
porque frequentemente confundido com automação de teste.
\end{conceito}

\section{Manutenção de software}

\begin{conceito}[Classificação pela causa da alteração, não pelo momento]
A ISO/IEC 14764 classifica o que se faz com o software \textbf{depois de
entregue}. O critério é sempre \textbf{a causa} da mudança:

\begin{center}
\begin{tabular}{@{}p{2.6cm}p{9.4cm}@{}}
\toprule
\textbf{Tipo} & \textbf{Causa da alteração} \\
\midrule
Corretiva & \textbf{corrigir defeito} já detectado (erro relatado, falha em
produção) \\
Adaptativa & acompanhar \textbf{mudança do ambiente} externo ao software:
sistema operacional, SGBD, plataforma, hardware, legislação \\
Perfectiva & \textbf{melhorar} o que já funciona: desempenho,
manutenibilidade, requisito novo pedido pelo usuário (sem que nada
estivesse “quebrado”) \\
Preventiva & corrigir uma \textbf{falha latente} antes que ela se
manifeste --- reduzir risco futuro identificado numa análise, não um
defeito já relatado \\
\bottomrule
\end{tabular}
\end{center}
\end{conceito}

\begin{exemplo}[Classificando um caso ambíguo]
“O time atualizou o driver de conexão com o banco de dados porque o
fabricante descontinuou a versão anterior, sem que houvesse defeito
relatado nem funcionalidade nova.” --- é \textbf{adaptativa}: a causa é o
\textbf{ambiente} (o fabricante mudou algo fora do controle do sistema).
Não é corretiva (não havia defeito a corrigir), não é perfectiva (não se
melhorou nada do ponto de vista do usuário) e não é preventiva (não havia
uma falha latente identificada e sim uma obsolescência externa).
\end{exemplo}

\begin{cuidado}
“Evolutiva” não é um dos quatro tipos da classificação --- é um rótulo
informal que às vezes aparece em prova como se fosse categoria da norma,
mas a ISO/IEC 14764 reconhece só as quatro da tabela acima.
\end{cuidado}

\section{Mensuração: Ponto de Função e Story Points}

\subsection{Ponto de Função (APF): medir de fora, independente da linguagem}

\begin{conceito}[O problema de contrato que a APF resolve]
A Análise de Pontos de Função resolve uma pergunta prática: \textit{como
pagar por software sem pagar por linha de código?} Medir por linhas de
código premia quem escreve mal e pune quem usa uma linguagem expressiva ---
a mesma funcionalidade pode sair em 500 linhas de uma linguagem e 80 de
outra, e não faz sentido a segunda valer menos.

A saída foi medir do \textbf{lado de fora}: conta-se o que o sistema
\textbf{faz para o usuário} (o que ele entrega em termos de negócio), não
como foi construído por dentro. Por isso a contagem parte de
\textbf{telas, relatórios e arquivos lógicos} --- artefatos visíveis ao
usuário --- e o resultado é \textbf{o mesmo} qualquer que seja a
linguagem, a equipe ou a ferramenta usada para construir. É essa
independência de tecnologia que torna a APF apta a comparar fornecedores
diferentes num mesmo contrato.

Contam-se cinco tipos de função, em dois grupos:

\begin{center}
\begin{tabular}{@{}p{2.6cm}p{9.6cm}@{}}
\toprule
\textbf{Funções de transação} & o que \textbf{atravessa} a fronteira do
sistema: \textbf{EE} (Entrada Externa), \textbf{SE} (Saída Externa),
\textbf{CE} (Consulta Externa) \\
\addlinespace
\textbf{Funções de dados} & o que o sistema \textbf{guarda ou lê}:
\textbf{ALI} (Arquivo Lógico Interno, mantido dentro do sistema),
\textbf{AIE} (Arquivo de Interface Externa, só lido, mantido por outro
sistema) \\
\bottomrule
\end{tabular}
\end{center}

Os dois critérios que resolvem quase toda classificação:

\begin{itemize}
\item \textbf{ALI x AIE se decide pela manutenção}, não pela leitura. Se o
sistema em análise \textbf{mantém} o dado (cria, atualiza, apaga), é ALI;
se ele só \textbf{consulta} um dado que outro sistema mantém, é AIE ---
mesmo que a leitura seja constante.
\item \textbf{SE x CE se decide pelo dado derivado}. Se a saída envolve
\textbf{cálculo, totalização ou lógica} que gera informação nova a partir
dos dados brutos, é Saída Externa; se é recuperação pura, sem processamento
que agregue valor, é Consulta Externa.
\end{itemize}
\end{conceito}

\begin{exemplo}[Classificando telas]
Uma tela de cadastro de cliente, que grava o registro no banco $\to$
\textbf{EE} (Entrada Externa: dado entra e é mantido). Um relatório que
soma o total de vendas do mês e mostra o percentual de crescimento em
relação ao mês anterior $\to$ \textbf{SE} (Saída Externa: há cálculo/
totalização). Uma tela que só busca e exibe o cadastro de um cliente pelo
CPF, sem processar nada $\to$ \textbf{CE} (Consulta Externa: recuperação
pura). A tabela de clientes, mantida por este mesmo sistema $\to$
\textbf{ALI}. Uma tabela de CEPs, consultada mas mantida por um sistema
externo dos Correios $\to$ \textbf{AIE}.
\end{exemplo}

\subsection{Ponto de Função x Story Points}

\begin{conceito}
\begin{center}
\begin{tabular}{@{}p{2.8cm}p{4.7cm}p{4.7cm}@{}}
\toprule
& \textbf{Ponto de Função (APF)} & \textbf{Story Points} \\
\midrule
Natureza & \textbf{objetiva}, padronizada por norma (IFPUG) & \textbf{relativa/subjetiva}, calibrada pelo time \\
Independe do time? & sim --- dois analistas certificados chegam ao mesmo
número & não --- o mesmo item vale pontos diferentes para times
diferentes \\
Bom para & contrato de escopo fechado, comparação entre projetos/fornecedores & estimativa ágil interna, planejamento de Sprint \\
Baseado em & funções do sistema (EE, SE, CE, ALI, AIE) & esforço/complexidade/incerteza percebidos pelo time \\
\bottomrule
\end{tabular}
\end{center}
A confusão mais comum inverte os dois: dizer que Story Point é objetivo e
comparável entre times, ou que serve a contrato de escopo fechado --- é
exatamente o oposto. Ponto de Função é o \textbf{padronizado e comparável};
Story Points é o \textbf{relativo ao time que o usa}, e por isso não deveria
ser usado para comparar a “produtividade” de dois times diferentes.
\end{conceito}

\section{DevOps e estratégias de entrega}

\begin{conceito}[CI, CD e DevOps]
\textbf{CI (Integração Contínua)}: integra e testa o código com frequência
(idealmente a cada \textit{commit}), pegando conflitos e regressões cedo.
\textbf{CD}: tem dois sentidos que convivem --- \textbf{Entrega Contínua}
(o software fica sempre em estado \textbf{pronto para ir} a produção,
mesmo que o “ir” seja manual) e \textbf{Implantação Contínua} (o “ir” para
produção também é automático). Na prática de prova, “fornecer rapidamente
uma nova versão em produção” é sempre \textbf{CD}, nunca CI --- CI para
antes de chegar a produção.

\textbf{DevOps} une desenvolvimento e operações num fluxo só, reduzindo o
atrito entre “quem constrói” e “quem opera”. \textbf{DevSecOps} injeta
verificações de segurança dentro desse mesmo pipeline, em vez de deixar a
segurança para uma revisão isolada no fim.
\end{conceito}

\begin{conceito}[Estratégias de implantação]
Formas de colocar uma nova versão em produção reduzindo o risco de queda:

\begin{itemize}
\item \textbf{Blue-green}: mantém \textbf{dois ambientes idênticos}
(azul = versão atual em produção, verde = nova versão) e troca o tráfego de
um para o outro de uma vez --- rollback é instantâneo, basta apontar de
volta para o azul.
\item \textbf{Canary}: libera a nova versão para uma \textbf{fatia pequena}
dos usuários primeiro, monitora, e só então expande gradualmente para
todos --- reduz o raio de impacto de um defeito não detectado em teste.
\item \textbf{Rolling update}: substitui as instâncias \textbf{aos poucos},
uma ou algumas de cada vez, sem nunca derrubar a capacidade total do
serviço.
\item \textbf{Feature toggle} (\textit{feature flag}): a nova
funcionalidade já está no código em produção, mas \textbf{desligada} por
uma chave de configuração --- liga-se para todos, para um grupo, ou nem se
liga, sem precisar de novo deploy.
\end{itemize}
\end{conceito}

\section{Design de software}

\begin{conceito}[Alto nível x baixo nível; arquitetura x design]
\textbf{Design de alto nível}: estrutura geral do sistema --- módulos e
como eles interagem, próximo do que se chama de \textbf{arquitetura}.
\textbf{Design de baixo nível}: detalhe interno de cada módulo --- funções,
métodos, algoritmos, estruturas de dados específicas.

\textbf{Arquitetura} trata de decisões \textbf{amplas e estruturais}
(quais componentes existem, como se comunicam); \textbf{design} (no sentido
estrito) trata de decisões \textbf{detalhadas} dentro de um componente já
definido. Não é verdade que “arquitetura só importa em projeto grande” ---
toda aplicação tem uma arquitetura, ainda que implícita e não documentada;
o risco de não decidir conscientemente é que ela emerja de forma
desorganizada.

Padrões de projeto (GoF) e notação UML têm tratamento completo no Capítulo
\ref{cap:padroes-uml-teoria} --- este capítulo cobre o \textbf{processo} em
volta do design; a técnica de design orientado a objetos propriamente dita
está lá.
\end{conceito}

\begin{regrapratica}
Ao ler uma questão que descreve uma decisão de projeto, pergunte: essa
decisão \textbf{muda a forma como os componentes se relacionam} (quantos
serviços existem, como se comunicam, onde os dados moram)? Se sim, é
arquitetura/alto nível. Ela só \textbf{muda o algoritmo/estrutura dentro de
um componente} sem afetar quem fala com quem? Se sim, é design/baixo
nível.
\end{regrapratica}
